\section{Discussion and Limitations}
\label{ch:discussion}

Five limitations bound what this thesis can claim. They are given here in full; seven further
caveats, each real but none of them load-bearing, are recorded in
Table~\ref{tab:further-caveats} of Appendix~\ref{app:supplementary}.

\subsection{The Baseline Is Calibrated, Not Validated}
\label{sec:lim-calibration}

The most important limitation in the thesis is that the baseline is calibrated rather than
validated. The income-shock probability of Submodel~1 cannot be estimated from the data, since
\ac{NIDS} Wave 5 enters this design as a single wave and within-household income volatility is not
identified from one wave, so it is fitted until baseline arrears reproduce pattern~2. Agreement
with the 2017 \ac{CCMR} on the fitted bands is consequently not evidence that the model is right;
only the bands left free, and the results obtained after \ac{BNPL} is injected with that parameter
held fixed, are genuine out-of-sample tests.

The fitted value is also implausibly high, reproducing observed arrears requiring a job-separation
rate roughly four times the one the labour force survey records. The reason is structural: this
model has exactly one thing that can go wrong, where real households fall behind through illness,
unexpected bills, rate rises and family change as well, so one channel is doing the work of
several. This is disclosed rather than corrected, since correcting it means adding shock channels
the data cannot parameterise.

\subsection{The Validation Targets Are Uneven}
\label{sec:lim-validation}

Only pattern~2 is both South African and vintage-matched to the 2017 population, and it is measured
on an account basis while the model's unit is the household. An account-level arrears share is not
a household-level distress rate, so the comparison asks only whether the profile is of the right
shape and scale; because pattern~2 is also the calibration target, that imprecision propagates into
the fitted shock probability.

Pattern~1 is not fully independent of the design. The want-driven borrowing path of Submodel~3 was
adopted partly because a purely shortfall-driven \ac{BNPL} agent would treat the product as a pure
substitute and could not reproduce the complementarity result under any parameterisation. The
direction of the effect is therefore partly built into the lender-choice rule; its magnitude is
not, and the pattern can still fail if the Regulation 23A gate crowds out traditional borrowing
sufficiently. Patterns 3 and 4 rest respectively on a United Kingdom credit-card model and on a
2025 survey compared across a period the counterfactual design otherwise forbids. All three are
accordingly reported as plausibility checks and not as tests.

\subsection{Behavioural Parameters Are Imported}
\label{sec:lim-transfer}

The behavioural micro-foundations of Section~\ref{sec:lit-behavioural} are drawn from United States
evidence and applied to South African borrowers: the minimum-payer share of 0.29, the present-bias
results, and the experiment that motivates the want-driven borrowing path. No developing-market
evidence of comparable quality was located for any of the three. The parameter that matters most is
swept from 0.20 to 0.40, and the mechanisms transferred are claims about general features of
decision-making under credit contracts rather than about United States institutions, but neither
argument establishes that the parameters are right for South Africa.

The coefficients of Equation~\ref{eq:peer} are fixed by no source at all. The $\beta = 0$ control
arm addresses the circularity this creates but not the plausibility problem, since it does not
identify which region of the $(q_{\text{base}}, \beta)$ surface corresponds to the actual market.
The available anchors are indirect: the 82.8\% banked ceiling, the 20\% of consumers naming
\ac{BNPL}, and the \ac{CFPB} stacking shares.

\subsection{One Rule Has No Anchor, and It Is the Largest Sensitivity}
\label{sec:lim-amount}

Submodel~4 governs how much a household asks for when it comes up short, and the literature search
located no study that settles it. Sensitivity analysis across three specifications is therefore
mandatory rather than optional, and this rule is the largest remaining bar on the tornado. Two of
the three agree to within the replicate noise floor and the whole range comes from the third and
most generous, so the sensitivity reflects one unconstrained variant rather than general
instability.

\subsection{The Model Cannot Produce a Cascade}
\label{sec:lim-structural}

Six exclusions were adopted by design: \textbf{no macroeconomic dynamics}, since any time-variation
would confound the injection effect with a macro effect; \textbf{no contagion}, there being one
aggregate lender with no balance sheet that can fail and no feedback from household default to
income; \textbf{no competition among traditional lenders}, which suppresses the lender-conduct
mechanism Hamill et al. show to matter; \textbf{no learning and no scarring}; \textbf{no peer
effects on consumption}, so peer influence acts only on \ac{BNPL} adoption; and \textbf{static
household composition}.

The second of these is not merely a caveat but the explanation of the thesis's central negative
result. With no channel by which one household's default affects another's income, what the model
produces is many correlated individual defaults rather than a cascade, and a threshold in the
population default rate requires feedback in the outcome itself. The absence of a cliff edge in
Section~\ref{sec:results-rq2} is therefore a property of a model built this way, and a successor
design with an impairable lender or a default-to-income channel could plausibly produce one. This
is stated here so that the negative result is read as bounded rather than as general.

\subsection{How Results Should Be Read}
\label{sec:lim-reading}

Results are directional and mechanistic: they describe what happens to a population of this kind
when a lender class with these properties is introduced, and are informative about the sign, the
rough magnitude and the mechanism of that change. They are not forecasts, and no number in
Section~\ref{ch:results} should enter a policy calculation without the qualifications above. What
survives all of them is the comparison the model was built to make, between two runs identical
except for the presence of a bureau-reporting obligation, which holds every uncalibrated parameter
fixed across both arms.

